\documentclass[../../main.tex]{subfiles}
\graphicspath{{../../image/}{../../figure/}}
\begin{document}

\section{注记}\label{section:ch04-注记}

在很大程度上，正是波破碎灾难为非线性双曲守恒律组理论定下基调。这一效应在第 4.2 小节中通过 Burgers 方程范式引入，但正如我们将在\cref{chapter:ch07-一维双曲平衡律组}中看到，它弥漫于单个空间变量中（真正）非线性守恒律组的全体。在多个空间维数中，波破碎的驱动仍然存在，但须与色散相竞争，后者可延迟、甚至完全阻止古典解在满足所谓\emph{零条件}的方程组中的崩溃。沿这一方向，从浩繁文献中，例如见 Christodoulou [1]、Klainerman [1]、Klainerman 与 Sideris [1]、Sideris [2,3,4]、Agemi [1]、Chae 与 Huh [1]、以及 Ta-tsien Li [1]。


特别地，对三维 Euler 方程组，波破碎与色散的竞争机制几乎势均力敌。因而，第 4.8 小节所注“波破碎最终占优”的证明需要非常精细的分析。在 John [2] 的开创性工作之后，古典解的崩溃由 Alinhac [1,2,3] 确立。然而，给出波破碎详细描述的确定性论述归功于：相对论 Euler 方程组方面 Christodoulou [2]，古典 Euler 方程组方面 Christodoulou 与 Miao [1]。这些证明非常技术化，占数百页篇幅。Holzegel、Klainerman、Speck 与 Wong [1] 的近期论文在本课题历史发展的语境下给出了非常可读的上述工作综述。Speck [1] 是一部专著，把结果推广到三维一般准线性波动方程组。

第 4.8 小节所勾勒的“质量囚禁是 Euler 方程组古典解崩溃的另一发动机”的证明，改编自 Sideris [1]。同一节中“气体质量体积在真空中以三次方阶膨胀、且此速率的增长无法维持除非界面处形成奇点”的论证取自 Serre [31]。在这一联系中，另见 Chemin [1]、Liu 与 Yang [1]、以及 Yang 与 Zhu [2]。关于在真空存在下用特征线方法证明 Euler 方程组古典解崩溃，见 Chae 与 Ha [1]。一类有趣的整体（“永恒”）古典解在 Serre [13,31] 与 Grassin 与 Serre [1] 中呈现于二维 Euler 方程组。另见 Chen 与 Young [3]。关于时间周期解，见 Georgiev 与 LeFloch [1]。

Cauchy 问题弱解的非唯一性以及发展选择判据的需要，对双曲守恒律理论构成另一挑战，其众多方面将在后续章节中广泛讨论。特别地，正如我们在历史导论以及第 II、III 章中所见，第 4.5 小节引入的熵可容性判据是 Clausius--Duhem 不等式的抽象，表达热力学第二定律；至少自二十世纪初以来它已被应用于具体情形。对此判据最早明确的抽象形式参考见于 Kružkov [1]，但其中心重要性在 Lax [4] 的开创性论文发表后才被认识，该文除其他外引入了当前语境下的“熵”一词。考虑到不等式\eqref{align:ch04-熵不等式}的方向与 Clausius--Duhem 不等式\cref{chapter:ch03-双曲平衡律方程组}（2.3.13）的方向相反，从连续物理立场看，使用“自由能”而非“熵”一词会更恰当。


正如我们在第 4.5 小节所见，熵的凸性--按 Lax [4] 所给概念定义是必要前提--为可容 $L^{\infty}$ 弱解诱导出一定程度的稳定性，但在连续物理中出现的方程组里并非总是满足。对更弱但仍充分的拟凸条件（已在第 4.5 小节注出），见 Morrey [1]、Dacorogna [1]、以及 Müller 与 Fonseca [1]。Šverák [1] 表明秩一凸一般不足以保证拟凸。多凸性由 Ball [1] 在弹性静力学语境中引入。我们将在第 5.4 小节回到这一概念。\cref{remark:ch04-多熵可容性}归功于 Gui-Qian Chen [9]。

关于极大耗散性思想的各种试验，见 Dafermos [32]、Demoulini、Stuart 与 Tzavaras [1]、Gangbo 与 Westdickenberg [1]、Westdickenberg [1]、Chiodaroli 与 Kreml [1]、以及 Feireisl [2]。我们将在第 9.7 小节、尽管在较窄范围内更详细地回到这一问题。

在后续章节中，我们将频繁遇到消失黏性方法，既作为构造解的方法，也作为识别可容激波的手段。正是为后一目的，该方法最初由 Rayleigh [4] 与 G.I. Taylor [1] 引入；见历史导论。

中间抛物--双曲型方程组理论的论述见于 Songmu Zheng [1] 与 Hsiao [3] 的专著，以及 Hsiao 与 Jiang [1] 的综述文章，读者可在其中找到广泛的文献列表。

关于弱解的黏性--毛细可容性条件，见 Slemrod [3] 与 LeFloch [5]。

关于初边值问题及相关文献的更详细讨论，读者应查阅第 V、VI 章。不等式\eqref{align:ch04-PBG}首先由 Bardos、Leroux 与 Nédélec [1] 对标量守恒律导出，后由 DuBois 与 LeFloch [1] 推广到单个空间变量中的方程组。正如我们将在\cref{chapter:ch06-标量守恒律L1理论}中看到，这些不等式在标量情形完全刻画了可容边界条件。


在 De Lellis 与 Székelyhidi [1,2] 以及 Buckmaster、De Lellis、Isett 与 Székelyhidi [1] 的开创性工作之后，本节第 4.8 小节所勾勒的 Euler 方程组奇异解理论在 Chiodaroli [1]、Chiodaroli 与 Kreml [1,2]、Chiodaroli、Feireisl 与 Kreml [1]、Feireisl [2]、Feireisl 与 Kreml [1]、Feireisl、Kreml 与 Vasseur [1]、Chiodaroli、De Lellis 与 Kreml [1]、以及 Villani [2] 中得到进一步发展。特别地，\cref{lemma:ch04-次解扰动}取自 Feireisl [2]。

\end{document}
